\section{Introduction}
\label{sec:intro}

The digital landscape is currently witnessing a convergence of text generation modalities.
Human communication is no longer confined to traditional keyboard typing; it increasingly involves speech-to-text (ASR) dictation and significant assistance from Large Language Models (LLMs).
Each of these sources leaves behind unique, systematic artifacts—phonetic errors from dictation, orthographic typos from keyboard proximity, and stylistic signatures from LLMs.
Understanding these 'modes' of origin is essential for applications ranging from authorship attribution to the development of robust, context-aware error correction systems.

{\bf Why does this research matter?}
As LLMs are deployed in more critical roles, the ability to interpret and verify the source of text becomes a cornerstone of digital trust.
Detecting whether a piece of text was dictated or typed can provide crucial context for forensic linguistics, while isolating AI-generated artifacts is vital for maintaining the integrity of human-centric datasets.

{\bf What is the gap in existing work?}
Current research in text source characterization is largely fragmented.
One body of work focuses heavily on the binary "Human vs. AI" detection problem \cite{ta2025faid}, while another addresses specific artifacts like ASR error correction \cite{havare2026lost} or stylometric signatures \cite{kumarage2023neural}.
However, there is a lack of research that generalizes these artifacts into a unified, interpretable 'mode' latent variable.
Existing models often 'gloss over' subtle phonetic or orthographic differences, failing to leverage the structured nature of these artifacts for source identification.

{\bf Our approach.}
In this work, we propose to treat the origin mode as a latent variable that can be both explicitly detected and implicitly isolated.
We construct a composite dataset that represents three primary modes: \aigen, \dictated (simulated via phonetic homophones), and \typed (simulated via keyboard proximity typos).
We then conduct two experiments: (1) a zero-shot classification task using \gptmini to evaluate explicit detectability, and (2) a latent embedding analysis using \minilm to measure the inherent clustering of modes in the embedding space (see \secref{sec:method}).


{\bf Quantitative preview.}
Our experiments reveal that LLMs possess a moderate capability for zero-shot mode detection, achieving \accuracy accuracy.
While the model perfectly identifies \aigen text (100\% recall), it completely fails to distinguish \dictated text from \typed or \aigen sources in zero-shot contexts.
Crucially, our embedding analysis yields a Silhouette Score of \silhouette, confirming that 'mode' is indeed a separable latent variable, with the primary axis of variance (\pcaone) corresponding to the presence and type of text artifacts.

Our main contributions are as follows:
\begin{itemize}[leftmargin=*,itemsep=0pt,topsep=0pt]
    \item We define and formalize the concept of 'mode' as a unified latent variable for text origin characterization.
    \item We conduct a zero-shot evaluation showing that while \aigen and \typed modes are detectable, phonetic artifacts in \dictated text present a significant challenge for current LLMs.
    \item We provide empirical evidence that origin mode is captured in standard text embeddings, showing moderate cluster separation.
    \item We release our composite dataset and experimental framework to facilitate further research into latent mode isolation.
\end{itemize}
